% This file is embedded in datatool-user.pdf version 3.0.1 2025-03-05
% Example 144 Multi Bar Chart With a Legend
% Label: "ex:barchartgrouplegend"
% arara: pdflatex
% arara: pdfcrop
\documentclass[12pt]{article}
\pagestyle{empty}
% sample CSV file:
\begin{filecontents}[noheader,overwrite]{studentmarks.csv}
Surname,Forename,StudentNo,Assign1,Assign2,Assign3
"Smith, Jr",John,102689,68,57,72
"Brown",Jane,102647,75,84,80
"Brown",Jane,102646,64,92,79
"Brown",Andy,103569,42,52,54
"Adams",Zoë,105987,52,48,57
"Brady",Roger,106872,68,60,62
"Verdon",Clare,104356,45,50,48
\end{filecontents}


\usepackage{databar}
% Load data from studentmarks.csv file:
\DTLsetup{store-datum,default-name=marks}
\DTLread{studentmarks.csv} 
\begin{document}
\DTLmultibarchart
{
 init={
  \renewcommand{\DTLbaratendtikz}{% 
    \node[at={(\DTLbarchartwidth,0pt)},anchor=south west] 
    {\begin{tabular}{l}
      {\DTLdobarcolor{1}\rule{\DTLbarwidth}{\DTLbarwidth}}
      Assignment 1\\
      {\DTLdobarcolor{2}\rule{\DTLbarwidth}{\DTLbarwidth}}
      Assignment 2\\
      {\DTLdobarcolor{3}\rule{\DTLbarwidth}{\DTLbarwidth}}
      Assignment 3
      \end{tabular}};
    }
  },
 variables={\assignI,\assignII,\assignIII},
 bar-width=12pt,
 group-label-align={center,top},
 bar-label={\xDTLinitials{\Forename}\xDTLinitials{\Surname}},
 y-ticks,axes=both
}
{marks}% database name
{
 \assignI=Assign1,
 \assignII=Assign2,
 \assignIII=Assign3,
 \Surname=Surname,
 \Forename=Forename
} 
\end{document}
